\section{Introducción}\label{cap:introduccion}

Los emprendimientos de servicios de belleza constituyen una actividad económica dinámica en las ciudades latinoamericanas y, en su mayoría, operan como microempresas con estructuras administrativas reducidas. En ellas, el registro de servicios, el control de la caja y el cálculo de las comisiones del personal se realizan habitualmente en hojas de cálculo elaboradas por la propia administradora. Si bien esta herramienta es accesible y flexible, la literatura especializada advierte que las hojas de cálculo de uso operativo presentan una incidencia recurrente de errores que rara vez se detectan a tiempo (Panko, 1998).


El presente trabajo aborda el análisis y la propuesta de solución a la gestión administrativa del salón de belleza \textit{Lush Beauty}, un emprendimiento en funcionamiento que lleva su control financiero en hojas de cálculo con una hoja por especialista. El diagnóstico preliminar, realizado sobre una planilla de ejemplo (agosto de 2026) que reproduce el formato y el tipo de datos que recibe el salón, evidenció servicios registrados en texto libre, combos sin descomponer, ausencia de método de pago y estructuras distintas entre hojas, condiciones que dificultan la conciliación de la caja y la liquidación semanal de sueldos.


Ante esta problemática, se desarrolló una investigación aplicada de tipo cuantitativo, con apoyo cualitativo, cuyo propósito central es diseñar y validar un sistema web móvil que centralice la caja y la nómina del salón, y evaluar su efecto sobre el tiempo de liquidación, la exactitud de los cálculos y la usabilidad percibida. El alcance geográfico corresponde al salón Lush Beauty, en la ciudad de Santa Cruz de la Sierra, Bolivia; el alcance temporal abarca desde la línea base retrospectiva (planilla de ejemplo de agosto de 2026 y registros previos del salón) hasta el periodo de prueba piloto prevista para enero de 2027; y el alcance temático se limita a los procesos de registro de servicios, cobro, propinas, deducciones y liquidación de comisiones.


Una feria tecnológica exige presentar un producto verificable, útil y diferenciable. Agregar módulos sin relación con el problema central podría elevar el costo de desarrollo y debilitar la evidencia. Por ello, la propuesta evoluciona el concepto desde “digitalizar la planilla” hacia “controlar, explicar y automatizar la operación financiera”, sin convertirlo en una plataforma generalista de reservas o comercio electrónico.


La innovación se organiza en dos capas complementarias. La primera garantiza un registro transaccional exacto, conciliable y auditable: servicios, cobros, descuentos, comisiones y liquidaciones. La segunda transforma esos datos en apoyo explícito a decisiones: indicadores de contribución, escenarios What-If, comparación de alternativas, advertencias sobre supuestos, recomendaciones justificadas y medición posterior de resultados. Estas funciones permiten distinguir entre describir el pasado, explicar una variación y decidir entre cursos de acción.


El documento diferencia: a) hallazgos del diagnóstico y datos ilustrativos; b) funciones decididas para un MVP de demostración; c) experimentos sujetos a evidencia y d) características reservadas a fases posteriores. El plan sirve como base de implementación, actualización del informe académico y presentación ante un jurado.


El documento se estructura en dieciséis capítulos. El primero corresponde a esta introducción; el segundo presenta el planteamiento del problema; el tercero y el cuarto exponen los objetivos y la justificación; el quinto desarrolla el marco teórico y tecnológico; el sexto contiene el benchmarking de sistemas similares y la diferenciación de la propuesta; el séptimo describe la propuesta funcional Lush Beauty Smart; el octavo presenta Lush Decision Studio, el modelo de soporte para decisiones; el noveno detalla el diseño y la arquitectura de software; el décimo y el undécimo formulan las hipótesis y su operacionalización; el duodécimo describe el marco metodológico y de evaluación; el decimotercero presenta el cronograma mediante diagramas de Gantt, el plan de demostración y los riesgos; el decimocuarto expone las conclusiones y recomendaciones; y los capítulos decimoquinto y decimosexto incluyen las referencias y los anexos.
